\chapter{Detailed Multi-Camera Results}
\label{app:metrics}

The four-camera study reports the multi-camera experiments at the level of
target groups and of the top five model families. This appendix gives the
complete tables behind them: the per-camera and per-subject validation
breakdowns, the full ten-model benchmarks at every stage, and the
thirteen-target metrics of the three approaches.

\section{Validation --- the four tests in full}

One table per test, in the order they attach to the pipeline.

\begin{table}[H]
\centering
\caption{Test 1 --- two-dimensional pose. Reprojection error (px) per camera
and subject.}
\label{tab:app:t1}
\small
\begin{tabular}{@{}lccccc@{}}
\toprule
\textbf{Subject} & \textbf{cam 50591643} & \textbf{cam 58860488} &
\textbf{cam 60457274} & \textbf{cam 65906101} & \textbf{Mean} \\
\midrule
\texttt{s03} & 9.8 & 6.5 & 8.0 & 7.8 & 8.0 \\
\texttt{s04} & 7.1 & 4.0 & 4.2 & 4.9 & 5.0 \\
\texttt{s05} & 9.2 & 5.6 & 9.9 & 9.5 & 8.6 \\
\texttt{s07} & 11.3 & 10.3 & 9.8 & 10.6 & 10.5 \\
\texttt{s08} & 17.9 & 11.9 & 14.1 & 10.2 & 13.5 \\
\texttt{s09} & 12.3 & 8.2 & 9.1 & 8.5 & 9.5 \\
\texttt{s10} & 8.5 & 5.1 & 5.4 & 5.8 & 6.2 \\
\texttt{s11} & 8.0 & 3.7 & 4.9 & 4.4 & 5.2 \\
\midrule
\textbf{Mean} & \textbf{10.5} & \textbf{6.9} & \textbf{8.2} & \textbf{7.7} &
\textbf{8.3} \\
\bottomrule
\end{tabular}
\end{table}

\begin{table}[H]
\centering
\caption{Test 2 --- three-dimensional reconstruction. Segment-direction error per
subject; the position RMSE is a corpus figure (\SI{27.8}{\milli\metre}).}
\label{tab:app:t2}
\small
\begin{tabular}{@{}lcc@{}}
\toprule
\textbf{Subject} & \textbf{$\text{dir}_{3D}$ (\si{\degree})} &
\textbf{threshold $<8$} \\
\midrule
\texttt{s03} & 5.46 & \ok \\
\texttt{s04} & 4.29 & \ok \\
\texttt{s05} & 6.46 & \ok \\
\texttt{s07} & 5.90 & \ok \\
\texttt{s08} & 7.13 & \ok \\
\texttt{s09} & 5.47 & \ok \\
\texttt{s10} & 5.70 & \ok \\
\texttt{s11} & 3.45 & \ok \\
\midrule
\textbf{Mean} & \textbf{5.48} & \ok \\
\bottomrule
\end{tabular}
\end{table}

\begin{table}[H]
\centering
\caption{Test 3 --- joint angles against Vicon. The shoulder range of motion
explains why the shoulder is not judged on correlation: in a controlled curl it
barely moves.}
\label{tab:mc:t3}
\small
\begin{tabular}{@{}lcccc@{}}
\toprule
\textbf{Subject} & \textbf{elbow $r$} & \textbf{elbow MAE (\si{\degree})} &
\textbf{shoulder MAE (\si{\degree})} & \textbf{shoulder ROM (\si{\degree})} \\
\midrule
\texttt{s03} & 0.991~\ok & 4.52~\ok & 1.05~\ok & 2.46 \\
\texttt{s04} & 0.990~\ok & 4.56~\ok & 1.00~\ok & 8.93 \\
\texttt{s05} & 0.993~\ok & 4.48~\ok & 1.14~\ok & 3.04 \\
\texttt{s07} & 0.996~\ok & 3.70~\ok & 0.88~\ok & 4.18 \\
\texttt{s08} & 0.992~\ok & 4.04~\ok & 1.39~\ok & 1.80 \\
\texttt{s09} & 0.991~\ok & 5.31~\warn & 2.01~\ok & 4.95 \\
\texttt{s10} & 0.995~\ok & 3.79~\ok & 1.03~\ok & 5.91 \\
\texttt{s11} & 0.994~\ok & 3.42~\ok & 0.78~\ok & 2.81 \\
\midrule
\textbf{Mean} & \textbf{0.993}~\ok & \textbf{4.23}~\ok & \textbf{1.16}~\ok & 4.26 \\
\bottomrule
\end{tabular}
\end{table}

\begin{table}[H]
\centering
\caption{Test 4 --- subject scaling. Segment lengths and the resulting
\textit{Arm26} scale factors, per subject.}
\label{tab:app:t4}
\small
\begin{tabular}{@{}lcccc@{}}
\toprule
\textbf{Subject} & \textbf{humerus (\si{\metre})} & \textbf{$\times$ humerus} &
\textbf{forearm (\si{\metre})} & \textbf{$\times$ forearm} \\
\midrule
\texttt{s03} & 0.287 & 0.986~\ok & 0.262 & 1.034~\ok \\
\texttt{s04} & 0.286 & 0.985~\ok & 0.260 & 1.026~\ok \\
\texttt{s05} & 0.270 & 0.927~\ok & 0.239 & 0.943~\ok \\
\texttt{s07} & 0.259 & 0.891~\ok & 0.227 & 0.893~\ok \\
\texttt{s08} & 0.234 & 0.805~\warn & 0.212 & 0.836~\warn \\
\texttt{s09} & 0.259 & 0.892~\ok & 0.231 & 0.912~\ok \\
\texttt{s10} & 0.308 & 1.061~\warn & 0.273 & 1.077~\warn \\
\texttt{s11} & 0.270 & 0.928~\ok & 0.242 & 0.954~\ok \\
\midrule
\textbf{Mean} & --- & \textbf{0.934}~\ok & --- & \textbf{0.959}~\ok \\
\bottomrule
\end{tabular}
\end{table}

\section{Complete model benchmarks}

The chapter tabulates the top five families at each stage. These are the ten.

\begin{table}[H]
\centering
\caption{Approach A --- all ten families, raw 11 inputs and with the 22 engineered
features. Fold-averaged mean R\textsuperscript{2}.}
\label{tab:app:A}
\small
\begin{tabular}{@{}lccr@{}}
\toprule
\textbf{Model} & \textbf{raw (11)} & \textbf{$+$ FE (22)} & \textbf{$\Delta$} \\
\midrule
LightGBM     & 0.920 & 0.950 & $+0.030$ \\
ExtraTrees   & 0.906 & 0.938 & $+0.031$ \\
LSTM         & 0.905 & 0.915 & $+0.010$ \\
RandomForest & 0.900 & 0.935 & $+0.035$ \\
Transformer  & 0.899 & 0.922 & $+0.022$ \\
MLP          & 0.856 & 0.908 & $+0.052$ \\
ANN          & 0.849 & 0.924 & $+0.074$ \\
CNN+LSTM     & 0.849 & 0.932 & $+0.083$ \\
1D-CNN       & 0.777 & 0.875 & $+0.097$ \\
Ridge        & 0.563 & 0.905 & $+0.343$ \\
\bottomrule
\end{tabular}
\end{table}

\begin{table}[H]
\centering
\caption{Approach B --- all ten families across the three stages.}
\label{tab:app:B}
\small
\begin{tabular}{@{}lccc@{}}
\toprule
\textbf{Model} & \textbf{raw 3D} & \textbf{$+$ features} &
\textbf{$+$ \SI{2}{\hertz} filter} \\
\midrule
ExtraTrees   & 0.756 & 0.861 & 0.897 \\
LightGBM     & 0.743 & 0.858 & 0.893 \\
Transformer  & 0.751 & 0.819 & 0.892 \\
RandomForest & 0.723 & 0.845 & 0.884 \\
Ridge        & 0.597 & 0.721 & 0.883 \\
MLP          & 0.553 & 0.739 & 0.876 \\
ANN          & 0.744 & 0.781 & 0.873 \\
LSTM         & 0.692 & 0.795 & 0.858 \\
CNN+LSTM     & 0.702 & 0.761 & 0.857 \\
1D-CNN       & 0.710 & 0.823 & 0.852 \\
\bottomrule
\end{tabular}
\end{table}

\begin{table}[H]
\centering
\caption{Approach C --- all ten families, raw 2D keypoints and after multi-view
fusion with feature engineering.}
\label{tab:app:C}
\small
\begin{tabular}{@{}lccr@{}}
\toprule
\textbf{Model} & \textbf{raw 2D} & \textbf{$+$ fusion \& features} &
\textbf{$\Delta$} \\
\midrule
LightGBM     & $-0.002$ & 0.845 & $+0.847$ \\
ExtraTrees   & $0.197$  & 0.839 & $+0.642$ \\
RandomForest & $0.062$  & 0.831 & $+0.769$ \\
Transformer  & $-0.015$ & 0.815 & $+0.829$ \\
LSTM         & $-0.172$ & 0.813 & $+0.985$ \\
CNN+LSTM     & $0.015$  & 0.810 & $+0.795$ \\
ANN          & $-0.420$ & 0.799 & $+1.219$ \\
MLP          & $-1.172$ & 0.793 & $+1.964$ \\
1D-CNN       & $-0.630$ & 0.765 & $+1.395$ \\
Ridge        & $-0.154$ & 0.642 & $+0.795$ \\
\bottomrule
\end{tabular}
\end{table}

\section{Thirteen-target metrics}

Pooled over the held-out predictions of all eight leave-one-subject-out folds.
Activation is dimensionless, normalised between 0 and 1.

\begin{table}[H]
\centering
\caption{Approach A --- per target.}
\label{tab:app:Afull}
\small
\begin{tabular}{@{}llccccc@{}}
\toprule
\textbf{Target} & \textbf{unit} & \textbf{R\textsuperscript{2}} & \textbf{RMSE} &
\textbf{MAE} & \textbf{Pearson $r$} & \textbf{NRMSE (\%)} \\
\midrule
\texttt{elbow\_moment} & \si{\newton\metre} & 0.921 & 0.850 & 0.456 & 0.961 & 3.63 \\
\addlinespace
\texttt{act\_BIClong}  & --- & 0.968 & 0.017 & 0.012 & 0.984 & 3.76 \\
\texttt{act\_BICshort} & --- & 0.960 & 0.010 & 0.007 & 0.981 & 4.38 \\
\texttt{act\_BRA}      & --- & 0.973 & 0.017 & 0.012 & 0.987 & 3.61 \\
\texttt{act\_BRD}      & --- & 0.962 & 0.008 & 0.006 & 0.981 & 4.08 \\
\addlinespace
\texttt{frc\_BIClong}  & \si{\newton} & 0.967 & 10.521 & 7.572 & 0.984 & 4.11 \\
\texttt{frc\_BICshort} & \si{\newton} & 0.965 &  3.596 & 2.401 & 0.983 & 3.69 \\
\texttt{frc\_BRA}      & \si{\newton} & 0.975 & 16.275 & 10.901 & 0.988 & 3.47 \\
\texttt{frc\_BRD}      & \si{\newton} & 0.962 &  1.681 & 1.198 & 0.981 & 2.98 \\
\addlinespace
\texttt{MF\_BIClong}   & \si{\percent} & 0.958 & 0.175 & 0.137 & 0.979 & 4.94 \\
\texttt{MF\_BICshort}  & \si{\percent} & 0.949 & 0.132 & 0.104 & 0.975 & 5.51 \\
\texttt{MF\_BRA}       & \si{\percent} & 0.894 & 0.153 & 0.100 & 0.947 & 8.77 \\
\texttt{MF\_BRD}       & \si{\percent} & 0.945 & 0.101 & 0.080 & 0.973 & 5.15 \\
\midrule
\textbf{Mean} & & \textbf{0.952} & & & & \\
\bottomrule
\end{tabular}
\end{table}

\begin{table}[H]
\centering
\caption{Approach B --- per target.}
\label{tab:appB}
\small
\begin{tabular}{@{}llccccc@{}}
\toprule
\textbf{Target} & \textbf{unit} & \textbf{R\textsuperscript{2}} & \textbf{RMSE} &
\textbf{MAE} & \textbf{Pearson $r$} & \textbf{NRMSE (\%)} \\
\midrule
\texttt{elbow\_moment} & \si{\newton\metre} & 0.819 & 1.288 & 0.832 & 0.905 & 5.51 \\
\addlinespace
\texttt{act\_BIClong}  & --- & 0.904 & 0.029 & 0.023 & 0.951 & 6.48 \\
\texttt{act\_BICshort} & --- & 0.941 & 0.013 & 0.009 & 0.970 & 5.32 \\
\texttt{act\_BRA}      & --- & 0.941 & 0.026 & 0.018 & 0.971 & 5.35 \\
\texttt{act\_BRD}      & --- & 0.899 & 0.013 & 0.010 & 0.948 & 6.63 \\
\addlinespace
\texttt{frc\_BIClong}  & \si{\newton} & 0.904 & 17.965 & 13.733 & 0.951 & 7.01 \\
\texttt{frc\_BICshort} & \si{\newton} & 0.934 &  4.928 &  3.577 & 0.967 & 5.06 \\
\texttt{frc\_BRA}      & \si{\newton} & 0.937 & 25.958 & 17.491 & 0.970 & 5.53 \\
\texttt{frc\_BRD}      & \si{\newton} & 0.892 &  2.837 &  2.163 & 0.945 & 5.03 \\
\addlinespace
\texttt{MF\_BIClong}   & \si{\percent} & 0.958 & 0.175 & 0.137 & 0.979 &  4.94 \\
\texttt{MF\_BICshort}  & \si{\percent} & 0.924 & 0.155 & 0.123 & 0.961 &  6.48 \\
\texttt{MF\_BRA}       & \si{\percent} & 0.702 & 0.210 & 0.128 & 0.852 & 12.04 \\
\texttt{MF\_BRD}       & \si{\percent} & 0.955 & 0.097 & 0.077 & 0.977 &  5.14 \\
\midrule
\textbf{Mean} & & \textbf{0.893} & & & & \\
\bottomrule
\end{tabular}
\end{table}

\begin{table}[H]
\centering
\caption{Approach C --- per target.}
\label{tab:appC}
\small
\begin{tabular}{@{}llccccc@{}}
\toprule
\textbf{Target} & \textbf{unit} & \textbf{R\textsuperscript{2}} & \textbf{RMSE} &
\textbf{MAE} & \textbf{Pearson $r$} & \textbf{NRMSE (\%)} \\
\midrule
\texttt{elbow\_moment} & \si{\newton\metre} & 0.778 & 1.424 & 0.972 & 0.884 & 6.09 \\
\addlinespace
\texttt{act\_BIClong}  & --- & 0.856 & 0.035 & 0.027 & 0.926 & 7.95 \\
\texttt{act\_BICshort} & --- & 0.917 & 0.015 & 0.011 & 0.958 & 6.31 \\
\texttt{act\_BRA}      & --- & 0.877 & 0.037 & 0.026 & 0.937 & 7.73 \\
\texttt{act\_BRD}      & --- & 0.847 & 0.016 & 0.013 & 0.921 & 8.17 \\
\addlinespace
\texttt{frc\_BIClong}  & \si{\newton} & 0.874 & 20.610 & 15.688 & 0.935 & 8.05 \\
\texttt{frc\_BICshort} & \si{\newton} & 0.917 &  5.540 &  3.991 & 0.958 & 5.69 \\
\texttt{frc\_BRA}      & \si{\newton} & 0.897 & 33.204 & 22.986 & 0.948 & 7.08 \\
\texttt{frc\_BRD}      & \si{\newton} & 0.841 &  3.434 &  2.683 & 0.918 & 6.09 \\
\addlinespace
\texttt{MF\_BIClong}   & \si{\percent} & 0.910 & 0.257 & 0.214 & 0.956 & 7.26 \\
\texttt{MF\_BICshort}  & \si{\percent} & 0.845 & 0.221 & 0.181 & 0.920 & 9.22 \\
\texttt{MF\_BRA}       & \si{\percent} & 0.716 & 0.245 & 0.170 & 0.860 & 11.24 \\
\texttt{MF\_BRD}       & \si{\percent} & 0.877 & 0.160 & 0.128 & 0.940 & 6.32 \\
\midrule
\textbf{Mean} & & \textbf{0.847} & & & & \\
\bottomrule
\end{tabular}
\end{table}

\section{Per-subject generalisation}

\begin{table}[H]
\centering
\caption{Mean R\textsuperscript{2} on each held-out subject, for the three
approaches.}
\label{tab:persubj}
\small
\begin{tabular}{@{}lccc@{}}
\toprule
\textbf{Held-out subject} & \textbf{A} & \textbf{B} & \textbf{C} \\
\midrule
\texttt{s03} & 0.924 & 0.881 & 0.830 \\
\texttt{s04} & 0.917 & 0.819 & 0.761 \\
\texttt{s05} & 0.982 & 0.949 & 0.910 \\
\texttt{s07} & 0.977 & 0.946 & 0.918 \\
\texttt{s08} & 0.974 & 0.916 & 0.884 \\
\texttt{s09} & 0.984 & 0.964 & 0.922 \\
\texttt{s10} & 0.965 & 0.929 & 0.924 \\
\texttt{s11} & 0.890 & 0.742 & 0.624 \\
\midrule
\textbf{Mean}  & \textbf{0.952} & \textbf{0.893} & \textbf{0.847} \\
\textbf{Worst} & 0.890 & 0.742 & 0.624 \\
\bottomrule
\end{tabular}
\end{table}

\section{Optuna searches --- the three approaches}

Each approach runs one independent \textbf{40-trial} TPE study per model family,
with the same objective throughout: mean R\textsuperscript{2} under a 3-fold
subject-grouped cross-validation. The tables below are the top ten of Approach~B,
the approach where the search changes the outcome. The trial retained by each
study is in bold.

\begin{table}[H]
\centering
\caption{Best value reached by each family, all three approaches. LightGBM wins
A and B; the MLP wins C by \num{0.007}.}
\label{tab:app:optbest}
\small
\begin{tabular}{@{}lccc@{}}
\toprule
\textbf{Family} & \textbf{A --- angles} & \textbf{B --- 3D joints} &
\textbf{C --- 2D keypoints} \\
\midrule
LightGBM     & \textbf{0.9457} & \textbf{0.8900} & 0.8275 \\
MLP          & 0.8992 & 0.8819 & \textbf{0.8349} \\
ExtraTrees   & 0.9057 & 0.8672 & 0.7834 \\
RandomForest & 0.8995 & 0.8612 & 0.7912 \\
\bottomrule
\end{tabular}
\end{table}

\begin{table}[H]
\centering
\caption{Approach B, LightGBM --- top ten of 40 trials on the Butterworth
features. Trial \#14 is retained.}
\label{tab:app:optBlgbm}
\scriptsize
\begin{tabular}{@{}lrrrrrrrr@{}}
\toprule
\textbf{trial} & \textbf{value} & \textbf{n\_est.} & \textbf{leaves} &
\textbf{lr} & \textbf{subsample} & \textbf{colsample} & \textbf{min\_child} &
\textbf{$\lambda$} \\
\midrule
\textbf{\#14} & \textbf{0.8900} & 330 & 19 & 0.1245 & 0.7886 & 0.9886 & 40 & 9.8716 \\
\#17 & 0.8892 & 260 & 23 & 0.0380 & 0.8500 & 0.9152 & 35 & 8.3809 \\
\#28 & 0.8891 & 241 & 23 & 0.0678 & 0.9532 & 0.9526 & 34 & 9.0993 \\
\#22 & 0.8888 & 276 & 19 & 0.0282 & 0.8621 & 0.8800 & 34 & 8.7120 \\
\#21 & 0.8888 & 251 & 18 & 0.0259 & 0.8632 & 0.8849 & 36 & 8.6951 \\
\#30 & 0.8887 & 331 & 35 & 0.0705 & 0.9472 & 0.9572 & 29 & 9.1394 \\
\#25 & 0.8887 & 329 & 15 & 0.0258 & 0.7894 & 0.8718 & 38 & 7.6932 \\
\#37 & 0.8886 & 227 & 30 & 0.0833 & 0.8910 & 0.8955 & 24 & 9.4238 \\
\#34 & 0.8886 & 222 & 23 & 0.0463 & 0.9067 & 0.9091 & 40 & 9.8753 \\
\#23 & 0.8885 & 298 & 20 & 0.0305 & 0.8395 & 0.9471 & 33 & 8.7449 \\
\bottomrule
\end{tabular}
\end{table}

Every one of the ten sits at $\lambda$ between \num{7.7} and \num{9.9}: the
search pushed the L2 penalty to the top of its range and kept it there, which is
the signature of an input the model must not be allowed to overfit.

\begin{table}[H]
\centering
\caption{Approach B, MLP --- top ten of 40 trials. The architecture is identical
in all ten; only the regularisation and the learning rate move.}
\label{tab:app:optBmlp}
\scriptsize
\begin{tabular}{@{}lrlrr@{}}
\toprule
\textbf{trial} & \textbf{value} & \textbf{hidden} & \textbf{alpha} &
\textbf{lr\_init} \\
\midrule
\textbf{\#29} & \textbf{0.8819} & $(128,128,64)$ & 0.0012 & 0.0040 \\
\#22 & 0.8815 & $(128,128,64)$ & 0.0000 & 0.0039 \\
\#32 & 0.8804 & $(128,128,64)$ & 0.0010 & 0.0041 \\
\#11 & 0.8802 & $(128,128,64)$ & 0.0001 & 0.0049 \\
\#38 & 0.8789 & $(128,128,64)$ & 0.0002 & 0.0033 \\
\#16 & 0.8775 & $(128,128,64)$ & 0.0000 & 0.0047 \\
\#37 & 0.8774 & $(128,128,64)$ & 0.0008 & 0.0049 \\
\#15 & 0.8771 & $(128,128,64)$ & 0.0000 & 0.0040 \\
\#21 & 0.8764 & $(128,128,64)$ & 0.0000 & 0.0044 \\
\#14 & 0.8763 & $(128,128,64)$ & 0.0000 & 0.0037 \\
\bottomrule
\end{tabular}
\end{table}

\begin{table}[H]
\centering
\caption{Approach B, the two forest ensembles --- top ten of 40 trials each.}
\label{tab:app:optBforest}
\scriptsize
\begin{tabular}{@{}lrrrr@{}lrrrr@{}}
\toprule
\multicolumn{5}{c}{\textbf{ExtraTrees}} & \multicolumn{5}{c}{\textbf{RandomForest}} \\
\cmidrule(lr){1-5}\cmidrule(lr){6-10}
\textbf{trial} & \textbf{value} & \textbf{n\_est.} & \textbf{depth} &
\textbf{max\_feat.} &
\textbf{trial} & \textbf{value} & \textbf{n\_est.} & \textbf{depth} &
\textbf{max\_feat.} \\
\midrule
\textbf{\#8}  & \textbf{0.8672} & 144 & 19 & 0.4113 & \textbf{\#35} & \textbf{0.8612} & 178 & 25 & 0.4060 \\
\#32 & 0.8672 & 163 & 19 & 0.6278 & \#38 & 0.8609 & 200 & 24 & 0.4003 \\
\#3  & 0.8667 & 160 & 27 & 0.4426 & \#12 & 0.8606 & 173 & 28 & 0.4016 \\
\#34 & 0.8666 & 186 & 17 & 0.6280 & \#8  & 0.8596 & 135 & 19 & 0.4113 \\
\#1  & 0.8665 & 139 & 21 & 0.6626 & \#31 & 0.8592 & 166 & 28 & 0.4123 \\
\#21 & 0.8665 & 139 & 19 & 0.6607 & \#11 & 0.8589 & 164 & 28 & 0.4159 \\
\#31 & 0.8664 & 138 & 19 & 0.6758 & \#21 & 0.8588 & 163 & 28 & 0.4036 \\
\#0  & 0.8664 & 157 & 23 & 0.7617 & \#19 & 0.8584 & 163 & 18 & 0.4154 \\
\#23 & 0.8664 & 145 & 22 & 0.6135 & \#25 & 0.8581 & 134 & 11 & 0.4089 \\
\#12 & 0.8663 & 189 & 15 & 0.4019 & \#32 & 0.8580 & 170 & 26 & 0.4383 \\
\bottomrule
\end{tabular}
\end{table}

RandomForest never leaves \texttt{max\_features} $\approx\num{0.40}$ in any of its
ten best trials, while ExtraTrees spreads from \num{0.40} to \num{0.76} for the
same score --- its extra randomisation already does what feature subsampling does
for the forest, so the parameter stops mattering.

\begin{table}[H]
\centering
\caption{Approach A, LightGBM --- top five of the 40-trial screening study.
Trial \#31 is its best. The shipped artefact comes from a separate, wider search
(trial \#11, 975 trees), reported in Section~\ref{sec:resI}.}
\label{tab:app:optunaA}
\scriptsize
\begin{tabular}{@{}lrrrrrrrr@{}}
\toprule
\textbf{trial} & \textbf{value} & \textbf{n\_est.} & \textbf{leaves} &
\textbf{lr} & \textbf{subsample} & \textbf{colsample} & \textbf{min\_child} &
\textbf{$\lambda$} \\
\midrule
\textbf{\#31} & \textbf{0.9457} & 203 & 15 & 0.1049 & 0.857 & 0.865 & 37 & 3.218 \\
\#34 & 0.9449 & 239 & 15 & 0.0463 & 0.833 & 0.858 & 32 & 3.087 \\
\#19 & 0.9448 & 287 & 32 & 0.0665 & 0.865 & 0.890 & 34 & 4.526 \\
\#32 & 0.9448 & 203 & 16 & 0.1537 & 0.857 & 0.902 & 33 & 2.901 \\
\#22 & 0.9446 & 300 & 19 & 0.0372 & 0.860 & 0.887 & 36 & 3.568 \\
\bottomrule
\end{tabular}
\end{table}
